

\vspace{-3mm}
\subsection{Architectural Paradigms}
\label{sec:background}

This subsection classifies the representative multi-agent LLM frameworks in Table~\ref{tab:github-stats-updated} into a small set of recurring \emph{architectural paradigms}. Each paradigm corresponds to a coherent class of framework-level design commitments that jointly fix orchestration $\mathcal{O}$, communication structure $\mathcal{C}$, agent abstraction, and environment coupling $\mathcal{E}$. These commitments constrain runtime coordination and execution, leading to distinct execution semantics, scalability, and performance trade-offs. We identify three such paradigms, graph-based, role-based, and Generative Agent-Based Modeling (GABM).


\vspace{-3mm}
\subsubsection{Graph-Based Frameworks.}
Graph-based frameworks encode orchestration as an explicit workflow modeled in a directed graph, where nodes represent agents or computational modules and edges define permissible control or data flows~\cite{langgraph2025, langflow2025, n8n2025}. Execution strictly follows this predefined structure, emphasizing deterministic control flow and transparent execution with coordination determined at design time and enforced at runtime.


\noindent\textbf{Definition 2.3 (Graph-Based Framework).}
A graph-based framework instantiates $\mathcal{O}$ as a directed graph $\mathcal{G}=(V,E)$, where $V=\{a_1,\ldots,a_n\}$ denotes agents or components and $E\subseteq V\times V$ defines allowable control or data flows. Communication $\mathcal{C}$ is restricted to information propagation along edges in $\mathcal{G}$, and no explicit environment state $\mathcal{E}$ is maintained by the framework.


\subsubsection{Role-Based Frameworks.}
Role-based frameworks organize multi-agent systems around textual role specifications that condition agent objectives, tool access, and interaction behavior~\cite{crewai2025,openaiagents2025,wu2024autogen,agno2025,OpenAgents}. Rather than prescribing fixed workflows, coordination emerges through role-conditioned reasoning, task delegation, and structured message exchange across agents. This paradigm offers more flexibility in control flow while still imposing structured interaction patterns such as hierarchical or sequential task organizations.



\noindent\textbf{Definition 2.4 (Role-Based Framework).}
A role-based framework instantiates a set of agents $\{a_i\}_{i=1}^{n}$ together with a role space $D$ and a mapping $R:\{a_i\}\rightarrow D$ assigning each agent a role. Orchestration $\mathcal{O}$ and communication $\mathcal{C}$ are governed by role-conditioned interactions, and no explicit environment state $\mathcal{E}$ is maintained.


\vspace{-2mm}
\subsubsection{GABM Frameworks.}
GABM frameworks treat multi-agent systems as simulations in which system behavior emerges from repeated agent--environment interactions rather than explicit orchestration between agents. Coordination is mediated indirectly through shared world state, where agents observe, act, and influence future system evolution. A representative system is \emph{Concordia}~\cite{vezhnevets2023generative}, which implements this paradigm through an explicit environment model.

\noindent\textbf{Definition 2.5 (Generative Agent-Based Modeling Framework).}
A GABM framework instantiates agents $\{a_i\}_{i=1}^{n}$ interacting with an environment state $\mathcal{E}$ governed by a transition function $T:\mathcal{E}\times\{a_i\}\rightarrow\mathcal{E}$. Orchestration $\mathcal{O}$ arises from environment-mediated interaction, and direct inter-agent communication $\mathcal{C}$ is not exposed.

